\documentclass[11pt,a4paper]{article}
\usepackage[T1]{fontenc}
\usepackage[margin=25mm]{geometry}
\usepackage{parskip,microtype}
\usepackage[hidelinks]{hyperref}
\pagestyle{empty}
\begin{document}
Dear Editor,

Please consider our manuscript, \textit{Model Selection for Energy-Aware Edge Intrusion Detection: A Controlled Measurement Study}, for publication as a regular paper in \textit{IEEE Transactions on Network and Service Management}.

The paper examines whether runtime detector selection provides benefits beyond directly deploying the detector selected by the controller. We compare static, tabular Q-learning, DQN, deadline-based, and confidence-cascade methods in ten paired physical blocks on a Raspberry Pi 4B, with varying arrival rates and attack prevalence. The study reports whole-device energy together with detection quality, completed work, and deadline outcomes.

Both learned controllers retain the lightweight detector and show no resolved energy benefit over its direct static deployment. The confidence cascade offers a measurable energy--detection trade-off relative to a static random forest. We also formalize the controller information boundary and a pairwise immediate-utility condition that clarifies how resource normalization influences action preferences. These results address the attribution of benefits in learned network-service control and the design of useful static and conditional baselines.

This submission follows the invitation to submit a new manuscript after the decision on TNSM-2026-10776, \textit{Energy-Aware Intrusion Detection System with Q-learning-Based Dynamic Scheduler on Lightweight IoT Devices}. The new paper expands the comparison methods, defines the state and service model, and replaces the earlier main evaluation with paired measurements and sensitivity analyses. A point-by-point response to the prior reviewers is enclosed. Code, configurations, and result-reproduction instructions accompany the study in the public repository.

Thank you for considering our manuscript.

Sincerely,\\
Nuonan Ouyang\\
On behalf of the authors\\
James Cook University\\
\href{mailto:nuonan.ouyang@my.jcu.edu.au}{nuonan.ouyang@my.jcu.edu.au}
\end{document}
